\section{模型评价与推广}
\subsection{模型优点}
\begin{enumerate}
\item \textbf{完整性有明确的退出依据。}各频道必须已清除或已认证不存在；源数未知时仍完成无遗漏搜索，达到 $16$ 个不同源后才使用数量上限缩减覆盖。两问独立测试均实现 $3000/3000$ 完成和 $100\%$ 源级清除率。
\item \textbf{几何保证与效率估计分别核验。}接收、覆盖和清除条件由有界误差集合与几何证书给出，候选评分只改变动作次序；不会以估计概率代替不存在证明。
\item \textbf{行动成本采用一致的全程口径。}移动、检测、切换、处置及无源确认均计入虚拟时间。最后清除与合法退出分开统计，避免把确认开销遗漏后误判效率。
\item \textbf{结论可追溯到同一模型与样本。}最终模型主结果、父模型规划对照、有利布局和官方交互单例分别标注身份，防止将不同算法或不同终点的数值拼接比较。
\end{enumerate}

\subsection{模型局限}
固定站点提供充分覆盖条件，并未证明站点数或总路线最优。几何外包和数值保护量可能延迟清除；进一步细分非凸可行域有望降低保守性。问题二的有限情景评分没有连续全局最优保证，且其完整选点器尚未直接接入多源部署策略。

最终 Q4 相对父模型的配对区间跨零，延长参数优化尚未形成可确认的显著收益。完整无源确认产生的长尾更值得关注，但任何删点、提前终止或路线变更都需重新验证退出依据。主动搜索得到的最短/最长布局只反映已检验候选，不能作为全局极值。

全部大样本性能来自自建研究分布；有限样本的 $100\%$ 完成不等于跨任意分布的统计保证。理论发现证书亦不自动保证在任意现实计算预算内完成：通信中断、数值边界和动作预算仍需工程核验。已归档的官方单例使用父模型，最终两问各三次正式测试及全程序运行时间仍待原始成绩确认。

\subsection{模型推广}
该方法可推广到有界观测误差下的移动传感器搜索：将传感器可见范围转化为覆盖条件，以平台的真实动作成本替换虚拟时间，再对合法动作排序。遇到障碍物或移动目标时，应分别引入可达路径或动态可行域传播，并重新证明覆盖与清除条件。
